\documentclass[10pt,a4paper]{amsart}
\usepackage{amsmath,amssymb,amsthm,graphicx}
% For commutative diagrams we use Xy-pic
\usepackage[all]{xy}

% We use 2cell for 2-commutative diagrams.
\xyoption{2cell}
% Theorem environments.
%
\theoremstyle{plain}
\newtheorem{theorem}[subsection]{Theorem}
\newtheorem{proposition}[subsection]{Proposition}
\newtheorem{lemma}[subsection]{Lemma}

\theoremstyle{definition}
\newtheorem{definition}[subsection]{Definition}
\newtheorem{example}[subsection]{Example}
\newtheorem{exercise}[subsection]{Exercise}
\newtheorem{situation}[subsection]{Situation}

\theoremstyle{remark}
\newtheorem{remark}[subsection]{Remark}
\newtheorem{remarks}[subsection]{Remarks}

\numberwithin{equation}{subsection}

% Macros
%
\def\lim{\mathop{\mathrm{lim}}\nolimits}
\def\colim{\mathop{\mathrm{colim}}\nolimits}
\def\Spec{\mathop{\mathrm{Spec}}}
\def\Hom{\mathop{\mathrm{Hom}}\nolimits}
\def\Ext{\mathop{\mathrm{Ext}}\nolimits}
\def\SheafHom{\mathop{\mathcal{H}\!\mathit{om}}\nolimits}
\def\SheafExt{\mathop{\mathcal{E}\!\mathit{xt}}\nolimits}
\def\Sch{\mathit{Sch}}
\def\Mor{\mathop{\mathrm{Mor}}\nolimits}
\def\Ob{\mathop{\mathrm{Ob}}\nolimits}
\def\Sh{\mathop{\mathit{Sh}}\nolimits}
\def\NL{\mathop{N\!L}\nolimits}
\def\CH{\mathop{\mathrm{CH}}\nolimits}
\def\proetale{{pro\text{-}\acute{e}tale}}
\def\etale{{\acute{e}tale}}
\def\QCoh{\mathit{QCoh}}
\def\Ker{\mathop{\mathrm{Ker}}}
\def\Im{\mathop{\mathrm{Im}}}
\def\Coker{\mathop{\mathrm{Coker}}}
\def\Coim{\mathop{\mathrm{Coim}}}

% Boxtimes
%
\DeclareMathSymbol{\boxtimes}{\mathbin}{AMSa}{"02}

%
% Macros for moduli stacks/spaces
%
\def\QCohstack{\mathcal{QC}\!\mathit{oh}}
\def\Cohstack{\mathcal{C}\!\mathit{oh}}
\def\Spacesstack{\mathcal{S}\!\mathit{paces}}
\def\Quotfunctor{\mathrm{Quot}}
\def\Hilbfunctor{\mathrm{Hilb}}
\def\Curvesstack{\mathcal{C}\!\mathit{urves}}
\def\Polarizedstack{\mathcal{P}\!\mathit{olarized}}
\def\Complexesstack{\mathcal{C}\!\mathit{omplexes}}
% \Pic is the operator that assigns to X its picard group, usage \Pic(X)
% \Picardstack_{X/B} denotes the Picard stack of X over B
% \Picardfunctor_{X/B} denotes the Picard functor of X over B
\def\Pic{\mathop{\mathrm{Pic}}\nolimits}
\def\Picardstack{\mathcal{P}\!\mathit{ic}}
\def\Picardfunctor{\mathrm{Pic}}
\def\Deformationcategory{\mathcal{D}\!\mathit{ef}}
\newenvironment{reference}{}{}
\newenvironment{slogan}{}{}
\newenvironment{history}{}{}
\usepackage{cleveref}
\crefname{originaldefinition}{original definition}{original definitions}
\Crefname{originaldefinition}{Original definition}{Original definitions}
\makeatletter
\newlabel{morphisms-definition-finite-type-point}{{02J1}{source}{}{sourcepointer-1}{}}
\newlabel{morphisms-definition-finite-type-point@cref}{{[originaldefinition][0][]02J1}{[1][0][]source}{}{}{}}
\makeatother
\crefname{originallemma}{original lemma}{original lemmas}
\Crefname{originallemma}{Original lemma}{Original lemmas}
\makeatletter
\newlabel{morphisms-lemma-identify-finite-type-points}{{02J2}{source}{}{sourcepointer-2}{}}
\newlabel{morphisms-lemma-identify-finite-type-points@cref}{{[originallemma][0][]02J2}{[1][0][]source}{}{}{}}
\makeatother
\crefname{originallemma}{original lemma}{original lemmas}
\Crefname{originallemma}{Original lemma}{Original lemmas}
\makeatletter
\newlabel{morphisms-lemma-dimension-fibre-at-a-point}{{02FX}{source}{}{sourcepointer-3}{}}
\newlabel{morphisms-lemma-dimension-fibre-at-a-point@cref}{{[originallemma][0][]02FX}{[1][0][]source}{}{}{}}
\makeatother
\crefname{originallemma}{original lemma}{original lemmas}
\Crefname{originallemma}{Original lemma}{Original lemmas}
\makeatletter
\newlabel{lemma-restrict-preserves-type}{{04MP}{source}{}{sourcepointer-4}{}}
\newlabel{lemma-restrict-preserves-type@cref}{{[originallemma][0][]04MP}{[1][0][]source}{}{}{}}
\makeatother
\crefname{originallemma}{original lemma}{original lemmas}
\Crefname{originallemma}{Original lemma}{Original lemmas}
\makeatletter
\newlabel{lemma-groupoid-on-field-dimension-equal-stabilizer}{{04MQ}{source}{}{sourcepointer-5}{}}
\newlabel{lemma-groupoid-on-field-dimension-equal-stabilizer@cref}{{[originallemma][0][]04MQ}{[1][0][]source}{}{}{}}
\makeatother
\crefname{originallemma}{original lemma}{original lemmas}
\Crefname{originallemma}{Original lemma}{Original lemmas}
\makeatletter
\newlabel{lemma-groupoid-on-field-locally-finite-type-dimension}{{04LU}{source}{}{sourcepointer-6}{}}
\newlabel{lemma-groupoid-on-field-locally-finite-type-dimension@cref}{{[originallemma][0][]04LU}{[1][0][]source}{}{}{}}
\makeatother
\crefname{originallemma}{original lemma}{original lemmas}
\Crefname{originallemma}{Original lemma}{Original lemmas}
\makeatletter
\newlabel{groupoids-lemma-restrict-stabilizer}{{04ML}{source}{}{sourcepointer-7}{}}
\newlabel{groupoids-lemma-restrict-stabilizer@cref}{{[originallemma][0][]04ML}{[1][0][]source}{}{}{}}
\makeatother
\crefname{originallemma}{original lemma}{original lemmas}
\Crefname{originallemma}{Original lemma}{Original lemmas}
\makeatletter
\newlabel{algebra-lemma-CM-dim-formula}{{00NA}{source}{}{sourcepointer-8}{}}
\newlabel{algebra-lemma-CM-dim-formula@cref}{{[originallemma][0][]00NA}{[1][0][]source}{}{}{}}
\makeatother
\crefname{originallemma}{original lemma}{original lemmas}
\Crefname{originallemma}{Original lemma}{Original lemmas}
\makeatletter
\newlabel{algebra-lemma-silly}{{00DS}{source}{}{sourcepointer-9}{}}
\newlabel{algebra-lemma-silly@cref}{{[originallemma][0][]00DS}{[1][0][]source}{}{}{}}
\makeatother
\crefname{originallemma}{original lemma}{original lemmas}
\Crefname{originallemma}{Original lemma}{Original lemmas}
\makeatletter
\newlabel{algebra-lemma-one-equation}{{00KW}{source}{}{sourcepointer-10}{}}
\newlabel{algebra-lemma-one-equation@cref}{{[originallemma][0][]00KW}{[1][0][]source}{}{}{}}
\makeatother
\crefname{originallemma}{original lemma}{original lemmas}
\Crefname{originallemma}{Original lemma}{Original lemmas}
\makeatletter
\newlabel{algebra-lemma-reformulate-CM}{{02JN}{source}{}{sourcepointer-11}{}}
\newlabel{algebra-lemma-reformulate-CM@cref}{{[originallemma][0][]02JN}{[1][0][]source}{}{}{}}
\makeatother
\crefname{originallemma}{original lemma}{original lemmas}
\Crefname{originallemma}{Original lemma}{Original lemmas}
\makeatletter
\newlabel{algebra-lemma-grothendieck-general}{{046Z}{source}{}{sourcepointer-12}{}}
\newlabel{algebra-lemma-grothendieck-general@cref}{{[originallemma][0][]046Z}{[1][0][]source}{}{}{}}
\makeatother
\crefname{originalequation}{original equation}{original equations}
\Crefname{originalequation}{Original equation}{Original equations}
\makeatletter
\newlabel{equation-restriction}{{04MN}{source}{}{sourcepointer-13}{}}
\newlabel{equation-restriction@cref}{{[originalequation][0][]04MN}{[1][0][]source}{}{}{}}
\makeatother
\title{Cohen–Macaulay groupoid slicing lowers orbit-fibre dimension by one}
\author{The Stacks Project Authors}
\date{stacks-project:more-groupoids.tex:local-snapshot-20261011; GNU FDL 1.2 or later}
\begin{document}
\maketitle
\noindent\textit{Editorial note.} In the displayed intermediate dimension quotient the author writes f; in the target local ring this means t(f), as explicitly used in the subsequent construction. The missing closing parenthesis and occasional i/j minimal-prime index slip are typographical and do not change the stated construction.\par\medskip
\noindent\textit{Allowed standard prerequisites (editorial).}\begin{itemize}
\item Explicit prerequisite boundary. Groupoid restriction and stabilizer base change, dimension of locally finite-type groupoids over a field, prime avoidance, dimension formula for Cohen–Macaulay local rings, and the standard flatness criterion for a fibrewise nonzerodivisor.
\end{itemize}\medskip
\ifdefined\hypertarget\hypertarget{sourcepointer-1}{}\fi
\noindent\textit{Source reference.} Original definition 02J1: focused-seeds/AGP/sources/stacks-snapshot-20261011/tags.txt, source lines 3320--3320. Stable Stacks tag for an expressly allowed foundational or notational input; the counted construction remains the two concrete core obligations. The retained local core is not replaced by this pointer. See standard\_prerequisites.\par\medskip
\ifdefined\hypertarget\hypertarget{sourcepointer-2}{}\fi
\noindent\textit{Source reference.} Original lemma 02J2: focused-seeds/AGP/sources/stacks-snapshot-20261011/tags.txt, source lines 3321--3321. Stable Stacks tag for an expressly allowed foundational or notational input; the counted construction remains the two concrete core obligations. The retained local core is not replaced by this pointer. See standard\_prerequisites.\par\medskip
\ifdefined\hypertarget\hypertarget{sourcepointer-3}{}\fi
\noindent\textit{Source reference.} Original lemma 02FX: focused-seeds/AGP/sources/stacks-snapshot-20261011/tags.txt, source lines 3208--3208. Stable Stacks tag for an expressly allowed foundational or notational input; the counted construction remains the two concrete core obligations. The retained local core is not replaced by this pointer. See standard\_prerequisites.\par\medskip
\ifdefined\hypertarget\hypertarget{sourcepointer-4}{}\fi
\noindent\textit{Source reference.} Original lemma 04MP: focused-seeds/AGP/sources/stacks-snapshot-20261011/tags.txt, source lines 6050--6050. Stable Stacks tag for an expressly allowed foundational or notational input; the counted construction remains the two concrete core obligations. The retained local core is not replaced by this pointer. See standard\_prerequisites.\par\medskip
\ifdefined\hypertarget\hypertarget{sourcepointer-5}{}\fi
\noindent\textit{Source reference.} Original lemma 04MQ: focused-seeds/AGP/sources/stacks-snapshot-20261011/tags.txt, source lines 6051--6051. Stable Stacks tag for an expressly allowed foundational or notational input; the counted construction remains the two concrete core obligations. The retained local core is not replaced by this pointer. See standard\_prerequisites.\par\medskip
\ifdefined\hypertarget\hypertarget{sourcepointer-6}{}\fi
\noindent\textit{Source reference.} Original lemma 04LU: focused-seeds/AGP/sources/stacks-snapshot-20261011/tags.txt, source lines 6020--6020. Stable Stacks tag for an expressly allowed foundational or notational input; the counted construction remains the two concrete core obligations. The retained local core is not replaced by this pointer. See standard\_prerequisites.\par\medskip
\ifdefined\hypertarget\hypertarget{sourcepointer-7}{}\fi
\noindent\textit{Source reference.} Original lemma 04ML: focused-seeds/AGP/sources/stacks-snapshot-20261011/tags.txt, source lines 6047--6047. Stable Stacks tag for an expressly allowed foundational or notational input; the counted construction remains the two concrete core obligations. The retained local core is not replaced by this pointer. See standard\_prerequisites.\par\medskip
\ifdefined\hypertarget\hypertarget{sourcepointer-8}{}\fi
\noindent\textit{Source reference.} Original lemma 00NA: focused-seeds/AGP/sources/stacks-snapshot-20261011/tags.txt, source lines 852--852. Stable Stacks tag for an expressly allowed foundational or notational input; the counted construction remains the two concrete core obligations. The retained local core is not replaced by this pointer. See standard\_prerequisites.\par\medskip
\ifdefined\hypertarget\hypertarget{sourcepointer-9}{}\fi
\noindent\textit{Source reference.} Original lemma 00DS: focused-seeds/AGP/sources/stacks-snapshot-20261011/tags.txt, source lines 508--508. Stable Stacks tag for an expressly allowed foundational or notational input; the counted construction remains the two concrete core obligations. The retained local core is not replaced by this pointer. See standard\_prerequisites.\par\medskip
\ifdefined\hypertarget\hypertarget{sourcepointer-10}{}\fi
\noindent\textit{Source reference.} Original lemma 00KW: focused-seeds/AGP/sources/stacks-snapshot-20261011/tags.txt, source lines 766--766. Stable Stacks tag for an expressly allowed foundational or notational input; the counted construction remains the two concrete core obligations. The retained local core is not replaced by this pointer. See standard\_prerequisites.\par\medskip
\ifdefined\hypertarget\hypertarget{sourcepointer-11}{}\fi
\noindent\textit{Source reference.} Original lemma 02JN: focused-seeds/AGP/sources/stacks-snapshot-20261011/tags.txt, source lines 3342--3342. Stable Stacks tag for an expressly allowed foundational or notational input; the counted construction remains the two concrete core obligations. The retained local core is not replaced by this pointer. See standard\_prerequisites.\par\medskip
\ifdefined\hypertarget\hypertarget{sourcepointer-12}{}\fi
\noindent\textit{Source reference.} Original lemma 046Z: focused-seeds/AGP/sources/stacks-snapshot-20261011/tags.txt, source lines 5490--5490. Stable Stacks tag for an expressly allowed foundational or notational input; the counted construction remains the two concrete core obligations. The retained local core is not replaced by this pointer. See standard\_prerequisites.\par\medskip
\ifdefined\hypertarget\hypertarget{sourcepointer-13}{}\fi
\noindent\textit{Source reference.} Original equation 04MN: focused-seeds/AGP/sources/stacks-snapshot-20261011/tags.txt, source lines 6049--6049. Stable Stacks tag for an expressly allowed foundational or notational input; the counted construction remains the two concrete core obligations. The retained local core is not replaced by this pointer. See standard\_prerequisites.\par\medskip
\begin{lemma}
\label{lemma-slice}
Let $S$ be a scheme.
Let $(U, R, s, t, c, e, i)$ be a groupoid scheme over $S$.
Let $G \to U$ be the stabilizer group scheme.
Assume $s$ and $t$ are Cohen-Macaulay and locally of finite presentation.
Let $u \in U$ be a finite type point of the scheme $U$, see
Morphisms, Definition \ref{morphisms-definition-finite-type-point}.
With notation as in
Situation \ref{situation-slice},
set
$$
d_1 = \dim(G_u), \quad
d_2 = \dim_{e(u)}(F_u).
$$
If $d_2 > d_1$, then there exist an affine scheme $U'$
and a morphism $g : U' \to U$ such that (with notation as in
Situation \ref{situation-slice})
\begin{enumerate}
\item $g$ is an immersion
\item $u \in U'$,
\item $g$ is locally of finite presentation,
\item the morphism $h : U' \times_{g, U, t} R \longrightarrow U$
is Cohen-Macaulay at $(u, e(u))$, and
\item we have $\dim_{e'(u)}(F'_u) = d_2 - 1$.
\end{enumerate}
\end{lemma}

\begin{situation}
\label{situation-slice}
Let $S$ be a scheme.
Let $(U, R, s, t, c)$ be a groupoid scheme over $S$.
Let $g : U' \to U$ be a morphism of schemes.
Let $u \in U$ be a point, and let $u' \in U'$ be a point such that
$g(u') = u$. Given these data, denote $(U', R', s', t', c')$
the restriction of $(U, R, s, t, c)$ via the morphism $g$.
Denote $G \to U$ the stabilizer group scheme of $R$, which
is a locally closed subscheme of $R$.
Denote $h$ the composition
$$
h = s \circ \text{pr}_1 : U' \times_{g, U, t} R \longrightarrow U.
$$
Denote $F_u = s^{-1}(u)$ (scheme theoretic fibre), and $G_u$ the
scheme theoretic fibre of $G$ over $u$.
Similarly for $R'$ we denote $F'_{u'} = (s')^{-1}(u')$.
Because $g(u') = u$ we have
$$
F'_{u'} = h^{-1}(u) \times_{\Spec(\kappa(u))} \Spec(\kappa(u')).
$$
The point $e(u) \in R$ may be viewed as a point on $G_u$ and $F_u$ also, and
$e'(u')$ is a point of $R'$ (resp.\ $G'_{u'}$, resp.\ $F'_{u'}$) which maps
to $e(u)$ in $R$ (resp.\ $G_u$, resp.\ $F_u$).
\end{situation}

\begin{proof}
Let $\Spec(A) \subset U$ be an affine neighbourhood of $u$
such that $u$ corresponds to a closed point of $U$, see
Morphisms, Lemma \ref{morphisms-lemma-identify-finite-type-points}.
Let $\Spec(B) \subset R$ be an affine neighbourhood of $e(u)$
which maps via $j$ into the open
$\Spec(A) \times_S \Spec(A) \subset U \times_S U$.
Let $\mathfrak m \subset A$ be the maximal ideal corresponding to $u$.
Let $\mathfrak q \subset B$ be the prime ideal corresponding to $e(u)$.
Pictures:
$$
\vcenter{
\xymatrix{
B & A \ar[l]^s \\
A \ar[u]^t
}
}
\quad\text{and}\quad
\vcenter{
\xymatrix{
B_{\mathfrak q} & A_{\mathfrak m} \ar[l]^s \\
A_{\mathfrak m} \ar[u]^t
}
}
$$
Note that the two induced maps
$s, t : \kappa(\mathfrak m) \to \kappa(\mathfrak q)$
are equal and isomorphisms as $s \circ e = t \circ e = \text{id}_U$.
In particular we see that $\mathfrak q$
is a maximal ideal as well. The ring maps $s, t : A \to B$ are
of finite presentation and flat. By assumption the ring
$$
\mathcal{O}_{F_u, e(u)} = B_{\mathfrak q}/s(\mathfrak m)B_{\mathfrak q}
$$
is Cohen-Macaulay of dimension $d_2$. The equality of dimension holds by
Morphisms, Lemma \ref{morphisms-lemma-dimension-fibre-at-a-point}.

\medskip\noindent
Let $R''$ be the restriction of $R$ to $u = \Spec(\kappa(u))$
via the morphism $\Spec(\kappa(u)) \to U$.
As $u \to U$ is locally of finite type,
we see that $(\Spec(\kappa(u)), R'', s'', t'', c'')$
is a groupoid scheme with $s'', t''$ locally of finite type, see
Lemma \ref{lemma-restrict-preserves-type}.
By
Lemma \ref{lemma-groupoid-on-field-dimension-equal-stabilizer}
this implies that $\dim(G'') = \dim(R'')$. We also have
$\dim(R'') = \dim_{e''}(R'') = \dim(\mathcal{O}_{R'', e''})$, see
Lemma \ref{lemma-groupoid-on-field-locally-finite-type-dimension}.
By
Groupoids, Lemma \ref{groupoids-lemma-restrict-stabilizer}
we have $G'' = G_u$. Hence we conclude that
$\dim(\mathcal{O}_{R'', e''}) = d_1$.

\medskip\noindent
As a scheme $R''$ is
$$
R'' =
R \times_{(U \times_S U)}
\Big(
\Spec(\kappa(\mathfrak m)) \times_S \Spec(\kappa(\mathfrak m))
\Big)
$$
Hence an affine open neighbourhood of $e''$ is the spectrum of the ring
$$
B \otimes_{(A \otimes A)} (\kappa(\mathfrak m) \otimes \kappa(\mathfrak m))
=
B/s(\mathfrak m)B + t(\mathfrak m)B
$$
We conclude that
$$
\mathcal{O}_{R'', e''} =
B_{\mathfrak q}/s(\mathfrak m)B_{\mathfrak q} + t(\mathfrak m)B_{\mathfrak q}
$$
and so now we know that this ring has dimension $d_1$.

\medskip\noindent
We claim this implies we can find
an element $f \in \mathfrak m$ such that
$$
\dim(B_{\mathfrak q}/(s(\mathfrak m)B_{\mathfrak q} + fB_{\mathfrak q}) < d_2
$$
Namely, suppose $\mathfrak n_j \supset s(\mathfrak m)B_{\mathfrak q}$,
$j = 1, \ldots, m$ correspond to the minimal primes of the local ring
$B_{\mathfrak q}/s(\mathfrak m)B_{\mathfrak q}$.
There are finitely many as this ring is Noetherian (since it is essentially
of finite type over a field -- but also because a Cohen-Macaulay ring is
Noetherian). By the Cohen-Macaulay condition we have
$\dim(B_{\mathfrak q}/\mathfrak n_j) = d_2$, for example by
Algebra, Lemma \ref{algebra-lemma-CM-dim-formula}.
Note that
$\dim(B_{\mathfrak q}/(\mathfrak n_j + t(\mathfrak m)B_{\mathfrak q}))
\leq d_1$
as it is a quotient of the ring
$\mathcal{O}_{R'', e''} =
B_{\mathfrak q}/s(\mathfrak m)B_{\mathfrak q} + t(\mathfrak m)B_{\mathfrak q}$
which has dimension $d_1$. As $d_1 < d_2$ this implies that
$\mathfrak m \not \subset t^{-1}(\mathfrak n_i)$.
By prime avoidance, see
Algebra, Lemma \ref{algebra-lemma-silly},
we can find $f \in \mathfrak m$ with $t(f) \not \in \mathfrak n_j$ for
$j = 1, \ldots, m$. For this choice of $f$ we have
the displayed inequality above, see
Algebra, Lemma \ref{algebra-lemma-one-equation}.

\medskip\noindent
Set $A' = A/fA$ and $U' = \Spec(A')$. Then it is clear that
$U' \to U$ is an immersion, locally of finite presentation
and that $u \in U'$. Thus (1), (2) and (3) of the lemma hold.
The morphism
$$
U' \times_{g, U, t} R \longrightarrow U
$$
factors through $\Spec(A)$ and corresponds to the ring map
$$
\xymatrix{
B/t(f)B \ar@{=}[r] & A/(f) \otimes_{A, t} B & A \ar[l]_-s
}
$$
Now, we see $t(f)$ is not a zerodivisor on
$B_{\mathfrak q}/s(\mathfrak m)B_{\mathfrak q}$ as this is a
Cohen-Macaulay ring of positive dimension and $f$ is not contained
in any minimal prime, see for example
Algebra, Lemma \ref{algebra-lemma-reformulate-CM}.
Hence by
Algebra, Lemma \ref{algebra-lemma-grothendieck-general}
we conclude that $s : A_{\mathfrak m} \to B_{\mathfrak q}/t(f)B_{\mathfrak q}$
is flat with fibre ring
$B_{\mathfrak q}/(s(\mathfrak m)B_{\mathfrak q} + t(f)B_{\mathfrak q})$
which is Cohen-Macaulay by
Algebra, Lemma \ref{algebra-lemma-reformulate-CM}
again. This implies part (4) of the lemma.
To see part (5) note that by Diagram (\ref{equation-restriction})
the fibre $F'_u$ is equal to the fibre of $h$ over $u$.
Hence
$\dim_{e'(u)}(F'_u) =
\dim(B_{\mathfrak q}/(s(\mathfrak m)B_{\mathfrak q} + t(f)B_{\mathfrak q}))$
by
Morphisms, Lemma \ref{morphisms-lemma-dimension-fibre-at-a-point}
and the dimension of this ring is $d_2 - 1$ by
Algebra, Lemma \ref{algebra-lemma-reformulate-CM}
once more. This proves the final assertion of the lemma and we win.
\end{proof}

\bibliographystyle{plain}
\bibliography{my}
\end{document}
